\chapter{Introduction}\label{ch:introduction}

\section{Motivation}\label{sec:motivation}

Colorectal cancer (CRC) is the third most common cancer globally, accounting for
approximately 10\% of all cancer diagnoses and the second leading cause of
cancer-related mortality~\cite{sung2021global}. Despite significant advances in
treatment, prognosis remains heavily dependent on the stage at which the disease
is detected. Tumours identified at an early, localised stage have five-year
survival rates exceeding 90\%, whereas metastatic disease drops below
15\%~\cite{siegel2020colorectal}. This stark contrast underscores the urgent
need for reliable, non-invasive early-detection strategies.

Traditional screening methods --- such as colonoscopy and faecal occult blood
tests --- suffer from low patient compliance, invasiveness, or limited
sensitivity. \emph{Liquid biopsy}, the analysis of circulating cell-free DNA
(cfDNA) fragments shed into the bloodstream by dying cells, has emerged as a
promising alternative. In cancer patients, a fraction of cfDNA originates from
tumour cells and carries somatic mutations as well as aberrant methylation
patterns. Detecting these tumour-derived fragments against the overwhelming
background of normal cfDNA is, however, a needle-in-a-haystack problem that
demands highly sensitive computational methods.

\section{Why Novelty Detection?}\label{sec:why-novelty}

A natural framing of this problem is \emph{novelty detection}: we learn a model
of ``normality'' from healthy donor DNA and then flag test reads that deviate
significantly from the learned distribution. This one-class formulation is
attractive because it requires no labelled tumour samples for training, side-stepping
the chronic scarcity of annotated cancer genomic data and avoiding biases towards
specific mutation signatures.

\section{Thesis Outline}\label{sec:outline}

The remainder of this thesis is organised as follows.

\begin{itemize}
    \item \textbf{Chapter~\ref{ch:background}} reviews the biological background
          of cfDNA and epigenetic modifications, surveys prior work on
          machine-learning approaches to cancer detection from sequencing data,
          and positions the mismatch kernel method within this landscape.
    \item \textbf{Chapter~\ref{ch:data}} describes the origin, pre-processing,
          and characteristics of the DNA sequence data used in this study.
    \item \textbf{Chapter~\ref{ch:kernel}} provides the mathematical foundations
          of kernel methods and develops the mismatch string kernel in detail,
          including algorithmic optimisations adopted in the implementation.
    \item \textbf{Chapter~\ref{ch:novelty}} introduces the novelty detection
          framework, focusing on the One-Class SVM and the patient-level
          Multiple Instance Learning aggregation.
    \item \textbf{Chapter~\ref{ch:experiments}} details the experimental
          protocol, including the Leave-One-Out cross-validation scheme and
          evaluation metrics.
    \item \textbf{Chapter~\ref{ch:results}} presents and discusses the
          experimental results.
    \item \textbf{Chapter~\ref{ch:conclusions}} summarises the contributions and
          outlines directions for future work.
    \item \textbf{Appendix~\ref{ch:implementation}} provides a practical guide
          to the software implementation, covering repository structure, key
          modules, and instructions for reproduction.
\end{itemize}